\documentclass[aspectratio=43]{beamer}
% aspectratio=169 ：16:9 寬螢幕（推薦）；改為 43 可換回 4:3

% ── 載入所有套件與全域設定 ──────────────────────────────────
\input{note-template/package_and_settings}


% ============================================================
%  【如何切換文件】
%
%  每份文件對應 Content/ 下的一個資料夾，
%  資料夾內固定包含：
%
%    0. 標題與作者.tex  ← 定義 \title / \author / \date
%    part0.tex          ← 標題頁 + 目錄
%    part0_all.tex      ← 依序 \input 所有 part 檔的索引
%    part1.tex, ...     ← 實際投影片內容（每個 part 為一組 frames）
%
% ── 資料夾結構 ────────────────────────────────────────────
%
%  Content/
%  ├── course/             ← 課堂講義
%  │   ├── {科目1}/
%  │   │   ├── note/note_1/      ← 上課投影片
%  │   │   ├── hw/hw_1/          ← 作業講解
%  │   │   ├── quiz/quiz_1/      ← 小考
%  │   │   └── exam/midterm|final/
%  │   └── {科目2}/   （結構同上）
%  │
%  ├── paper/              ← 論文報告（讀書會、課堂）
%  │   ├── {作者_年份1}/
%  │   └── {作者_年份2}/
%  │
%  ├── thesis/             ← 論文口試簡報
%  │   └── {論文口試}/
%  │
%  └── meetings/           ← 會議報告（按報告對象 → 日期）
%      └── {對象}/
%          └── {YYYYMMDD}/   ← 日期格式：YYYYMMDD
%
%  圖片放在各文件的 figures/ 子資料夾，引用時使用完整路徑：
%    \includegraphics[width=0.8\textwidth]
%      {Content/meetings/{對象}/{YYYYMMDD}/figures/my_fig}
%
% ── 切換步驟（只需改以下兩行）────────────────────────────
%
%  步驟一：將下方 \input 換成目標文件的「0. 標題與作者.tex」路徑
%  步驟二：將 \begin{document} 內的 \input 換成目標的「part0_all.tex」路徑
%
%  範例 A ── 會議報告（{對象}，{日期}）：
%    \input{Content/meetings/{對象}/{YYYYMMDD}/0. 標題與作者}
%    \input{Content/meetings/{對象}/{YYYYMMDD}/part0_all}
%
%  範例 B ── 論文報告（Babii et al. 2022）：
%    \input{Content/paper/{作者_年份1}/0. 標題與作者}
%    \input{Content/paper/{作者_年份1}/part0_all}
%
%  範例 C ── 碩士論文口試：
%    \input{Content/thesis/{論文口試}/0. 標題與作者}
%    \input{Content/thesis/{論文口試}/part0_all}
%
%  範例 D ── {科目}第 1 章講義：
%    \input{Content/course/{科目1}/note/note_1/0. 標題與作者}
%    \input{Content/course/{科目1}/note/note_1/part0_all}
% ============================================================

% ── 步驟一：修改此行以載入標題與作者資訊 ─────────────────
\input{Content/course/__SUBJECT__/note/note_1/0. 標題與作者}

\begin{document}

% ── 步驟二：修改此行以載入文件內容 ──────────────────────
\input{Content/course/__SUBJECT__/note/note_1/part0_all}

\end{document}
